% Complete PMCE proofs and source examples for ED26-P18-erotetic-0001.
\section{Question evocation: proofs}
\label{sec:annex-p18-question-evocation}

Use the syntax, classical valuations and exact rule sheet of
Chapter~\ref{sec:p18-question-evocation}. The constructions below supply
the source's soundness, normal-form and finite-premise completeness
results and all its displayed derived or admissible transformations
\cite[secs.~2--5]{Wisniewski2016QuestionEvocation}.

\begin{annexitem}{annex:p18.qev.semantic-profile}{thm:p18-qev-semantic-profile}{The semantic criterion and its boundaries}
\begin{proof}
For a fixed valuation, the finite disjunction
\(A_1\lor\cdots\lor A_n\) is true exactly when at least one
\(A_i\) is true. This follows by induction on the displayed
disjunction using the binary truth clause. Quantifying over valuations
of \(X\) gives the coverage equivalence. Properness is precisely
the additional collection of non-entailments in the statement.

If \(X\) has no satisfying valuation, \(X\models A_i\) for
every \(i\), so properness fails. If an answer \(A_i\) is a
tautology, every \(X\)-valuation satisfies it and the same condition
fails. If \(A_i\leftrightarrow C_i\) is a tautology for every
\(i\), then the two lists have the same truth values coordinate by
coordinate at every valuation. Their coverage and each individual
non-entailment therefore agree. The replacement must independently
satisfy the syntax of a question. Finally, permuting a list changes
neither existence of a true answer nor the collection of individual
non-entailments, proving permutation invariance.
\end{proof}
\end{annexitem}

\begin{annexitem}{annex:p18.qev.soundness}{thm:p18-qev-soundness}{Axiom and primitive-rule soundness}
\begin{proof}
For an axiom, each answer clause lacks a complementary pair. Set every
positive literal's variable to zero and every negative literal's
variable to one. These instructions are compatible precisely because
there is no complementary pair. Extend the assignment to all other
variables by zero. It falsifies that clause. Thus each answer has a
countervaluation. The disjunction of all answer clauses contains some
\(p\) and \(\neg p\); every valuation satisfies one of them and
therefore at least one answer. The axiom evokes its question.

For \(R_1\), suppose both premise questions are evoked by \(X\).
Take a valuation satisfying \(X\). If it makes a common answer true,
the conclusion is covered. Otherwise the two premise coverage
conditions force both \(B\) and \(C\) to be true, so \(B\land C\)
is true. No common answer is entailed, by either premise. If
\(X\models B\land C\), then \(X\models B\), contradicting
unsettledness of the first premise. Hence the new answer is not entailed.
The rule proviso ensures that the resulting expression is a question.

For \(R_2\), the assumed tautology \(B\leftrightarrow C\)
makes the changed answers agree at every valuation. Coverage is
unchanged, and \(X\models C\) would imply \(X\models B\).
All common non-entailments persist. Syntactic distinctness is secured
by the proviso, rather than inferred from semantic equivalence.

For \(R_3\), let the conditional question be evoked by \(X\).
Any valuation of \(X\cup\{B\}\) must make some \(B\to A_i\)
true; since it makes \(B\) true, it makes that \(A_i\) true.
For every \(i\), non-entailment of \(B\to A_i\) from \(X\)
supplies a valuation satisfying \(X\), making \(B\) true and
\(A_i\) false. This witnesses the required non-entailment from
\(X\cup\{B\}\). Fixed-prefix implication is syntactically
injective, so the unprefixed answers remain distinct.

For \(R_4\), equality of direct-answer sets leaves exactly the same
coverage and non-entailment conditions. These four arguments establish
rule preservation. Induction on the finite proof tree now proves
soundness: leaves are sound axioms and every internal node is sound
when its immediate predecessors are sound.
\end{proof}
\end{annexitem}

\begin{annexitem}{annex:p18.qev.classical-tools}{thm:p18-qev-classical-tools}{CNF existence, freshness and restoration}
\begin{proof}
First take a formula \(A\), and list its finitely many variables as
\(p_1,\ldots,p_r\). There is at least one because the language has
no constants. For each Boolean assignment \(w\) to these variables
that falsifies \(A\), form a clause \(D_w\) containing \(p_i\)
when \(w(p_i)=0\) and \(\neg p_i\) when \(w(p_i)=1\).
This clause is false exactly at \(w\). Consequently, the conjunction
of all such \(D_w\) is true exactly at assignments satisfying \(A\).
It is a finite nonempty CNF whenever \(A\) is not a tautology.
If \(A\) is a tautology, use the single clause
\(p_1\lor\neg p_1\). Extensions to the remaining variables do
not affect the formulas. This proves CNF existence without assuming
a separate normal-form algorithm.

A clause with complementary literals is a tautology by the classical
negation and disjunction clauses. Without a complementary pair the
falsifying assignment used in the soundness proof is compatible, so
the clause is not a tautology. This proves both directions.

For a clause \(D\), choose one of its literals \(L\). Repeatedly
append another occurrence of \(L\) by disjunction, with a fixed
parenthesization. Every resulting expression remains a clause
equivalent to \(D\), while its syntactic length strictly increases.
Only finitely many lengths occur in a finite forbidden collection;
choose a variant longer than all of them. Applying this construction
successively produces any finitely many required clause variants that
are pairwise distinct and avoid a specified finite set. The method
also works when several original clauses are equivalent or identical.

For the question-normalization claim, first choose CNFs
\(N_1,\ldots,N_n\) equivalent to the original answers. Choose
distinct variables \(z_1,\ldots,z_n\) absent from every original
answer and every \(N_i\). This is possible because the union of their
variable sets is finite and \(\mathrm{Var}\) is infinite. Put
\[
 C_i=N_i\land(z_i\lor\neg z_i).
\]
Each \(C_i\) is a CNF equivalent to \(A_i\). It contains
\(z_i\), which no original answer contains; hence it equals no
\(A_j\). It also cannot equal another \(C_j\): for \(j\ne i\),
the latter contains no \(z_i\). Thus the new list is a question
and its answer set is disjoint from the original answer set.
The semantic profile theorem gives preservation of evocation.

Finally suppose a derivation with antecedent \(X\) ends in a
question with answers \(C_1,\ldots,C_n\), equivalent coordinatewise
to \(A_1,\ldots,A_n\), and the two answer sets are disjoint.
Use \(R_4\) to place any unreplaced \(C_i\) last, then use
\(R_2\) to replace it by \(A_i\). At that stage all other answers
are either unreplaced \(C_j\), which cannot equal \(A_i\) by
disjointness, or already restored \(A_j\), which cannot equal
\(A_i\) by the original question syntax. Thus every \(R_2\)
proviso holds. After finitely many replacements use \(R_4\) to
restore the desired order. This proves the claimed proof transformation.
\end{proof}
\end{annexitem}

\begin{annexitem}{annex:p18.qev.representative-induction}{thm:p18-qev-representative-induction}{A well-founded construction with syntactic markers}
\begin{proof}
We prove the assertion for all questions with a chosen representative
family of rank \(k\), by strong induction on the nonnegative integer
\(k\). The actual answer expressions need not be in CNF. This
distinction permits syntactic markers without increasing the induction
parameter of a recursive call.

If \(k=0\), each representative is one clause \(D_i\). Since
\(E(\varnothing,Q)\), each \(A_i\), and therefore each
\(D_i\), is nontautological. Their disjunction is a tautology by
coverage. Use the preceding lemma to replace the \(D_i\) by
equivalent clause variants \(D_i'\) that are pairwise syntactically
distinct and avoid all original \(A_j\). Each \(D_i'\) lacks a
complementary pair, and their disjunction is tautological, hence contains
a complementary pair. Thus \(\vdash?\{D_1',\ldots,D_n'\}\)
is an axiom. Restore the original answers using the preceding lemma's
\(R_2,R_4\) transformation. This proves the base case even if
the unmodified representative clauses were not syntactically distinct.

Let \(k>0\). Choose an answer \(A_j\) with representative
\(C_j=B_1\land\cdots\land B_m\), where \(m>1\) and all
\(B_i\) are clauses. This representative is nontautological because
\(A_j\) is. Some \(B_e\) must therefore be nontautological:
otherwise all its conjuncts, and the conjunction, would be tautologies.
Write \(D\) for the conjunction of the remaining \(m-1\) clauses,
retaining their order; it is nonempty. Let \(\vec A\) be the
original list with \(A_j\) removed, so \(\vec A\) is nonempty.

Choose distinct fresh variables \(u,v\) absent from all actual
answers and all chosen representatives, and put
\[
 U=B_e\land(u\lor\neg u),\qquad
 V=D\land(v\lor\neg v).
\]
Both \(?\{\vec A,U\}\) and \(?\{\vec A,V\}\) are questions:
each new last answer has a fresh variable not present in any common
answer. The unchanged answers are pairwise distinct. A valuation making
every common answer false must, by coverage of \(Q\), make \(A_j\)
and hence every \(B_i\) true. It therefore makes both \(U\)
and \(V\) true. This proves coverage for both new lists.

The common answers remain nontautological. Since \(U\) is equivalent
to the nontautological clause \(B_e\), the first new question is
evoked. Choose \(B_e\), without its marker, as its CNF representative;
keep all the other representatives. Its representative rank is
\(k-(m-1)<k\). The induction hypothesis yields a proof of
\(\vdash?\{\vec A,U\}\).

There are now two cases.

\emph{The residual \(D\) is nontautological.} Then \(V\) is
nontautological, so the second new question is evoked. Choose \(D\)
as its representative. Its rank is
\[
 k-(m-1)+(m-2)=k-1<k.
\]
The induction hypothesis supplies a proof of
\(\vdash?\{\vec A,V\}\). Apply \(R_1\) to the two proved
sequents, giving \(\vdash?\{\vec A,U\land V\}\).
This application is well formed: the conjunction contains the fresh
markers and therefore equals no common answer. The last answer is
classically equivalent to \(C_j\) and hence to \(A_j\).
Apply \(R_2\) to replace it by \(A_j\); the original question
syntax ensures that \(A_j\) equals no common answer. A final
\(R_4\) restores its original position.

\emph{The residual \(D\) is a tautology.} In this case \(V\)
is a tautology, so no induction call on its question is made. Instead
\(A_j\) is equivalent to \(B_e\land D\), hence to \(B_e\)
and to \(U\). Apply \(R_2\) directly to the already proved
\(\vdash?\{\vec A,U\}\), replacing \(U\) by \(A_j\),
then apply \(R_4\). The same original-distinctness argument
secures the replacement proviso.

All recursive calls have smaller representative rank. Each node uses
finitely many primitive-rule steps, and branching is at most binary.
Induction therefore constructs a finite PMCE proof in both cases.
\end{proof}
\end{annexitem}

\begin{annexitem}{annex:p18.qev.empty-completeness}{thm:p18-qev-empty-completeness}{From normal forms to arbitrary empty-antecedent questions}
\begin{proof}
For a normal-form question use each actual answer as its own CNF
representative. The preceding lemma then proves its sequent.

For arbitrary \(Q=?\{A_1,\ldots,A_n\}\) evoked by the empty
set, the classical-tools lemma constructs a question with pairwise
distinct CNF answers \(C_i\), equivalent to the \(A_i\) and
disjoint from them. The semantic profile gives its evocation. The
normal-form result supplies a proof of its sequent. Successively
restore the original answers with \(R_2,R_4\), using disjointness
to keep every intermediate question well formed. The resulting root
is \(\vdash Q\).
\end{proof}
\end{annexitem}

\begin{annexitem}{annex:p18.qev.conditionalization}{thm:p18-qev-conditionalization}{Conditionalization preserves the exact evocation criterion}
\begin{proof}
For a valuation satisfying \(X\), if \(B\) is false then every
\(B\to A_i\) is true. If \(B\) is true, one of these
implications is true exactly when one \(A_i\) is true. Thus the
conditional list covers all \(X\)-valuations exactly when the
original list covers all \(X\cup\{B\}\)-valuations.

For each \(i\), a countervaluation to
\(X\models B\to A_i\) is exactly a valuation satisfying
\(X\cup\{B\}\) and falsifying \(A_i\). Hence
\[
 X\not\models B\to A_i
 \quad\Longleftrightarrow\quad
 X\cup\{B\}\not\models A_i.
\]
Coverage and all individual non-entailments therefore agree in both
directions. This argument includes inconsistent premise sets: in that
case the relevant non-entailments fail, rather than an evocation being
obtained by vacuity.

In the fixed formula grammar, \(B\to A_i=B\to A_j\) as
expressions implies \(A_i=A_j\). Thus prefixing does not identify
distinct answers. Repeatedly apply the equivalence for an enumeration
\(B_1,\ldots,B_m\) of the finite set \(X\). It produces exactly
the nested prefixes in the statement. The empty enumeration changes
nothing and gives the stated convention for \(I_{\varnothing}\).
\end{proof}
\end{annexitem}

\begin{annexitem}{annex:p18.qev.completeness}{thm:p18-qev-completeness}{Finite-premise adequacy}
\begin{proof}
Soundness has already been proved. For completeness, assume \(E(X,Q)\).
If \(X=\varnothing\), use empty-antecedent completeness. Otherwise
write \(X=\{B_1,\ldots,B_m\}\) and
\(Q=?\{A_1,\ldots,A_n\}\). Conditionalization gives evocation
of \(?\{I_X(A_1),\ldots,I_X(A_n)\}\) by the empty set. The
empty-antecedent result provides its proof.

Apply \(R_3\) first with \(B_1\), leaving the common prefix
beginning at \(B_2\) on every answer and antecedent \(\{B_1\}\).
Continue with \(B_2\), and so on. After \(m\) applications,
the antecedent is exactly \(X\) and the question is exactly \(Q\).
At each stage the remaining common prefix is injective on answer
syntax; therefore there are still \(n>1\) distinct answers.
Every step is a primitive PMCE rule. Appending these finitely many
steps to the existing finite proof proves the desired sequent.
\end{proof}
\end{annexitem}

\begin{annexitem}{annex:p18.qev.derived-rules}{thm:p18-qev-derived-rules}{Moving the active answer position}
\begin{proof}
For \(R_2^*\), apply \(R_4\) to move the chosen answer to
the rightmost position, retaining a fixed order of the other answers.
Apply \(R_2\) there, then \(R_4\) to put the replacement in
the original position. The given distinctness proviso makes the middle
and final lists questions. The two permutations preserve their answer
sets, so they are legitimate \(R_4\) instances.

For \(R_1^*\), perform the same initial permutation on each of
the two premise questions, using the same order of common answers.
Apply \(R_1\) to their rightmost answers, then permute back.
The conjunction-distinctness proviso is exactly what the middle rule
and final list require. Transposing two answer positions is already
a single \(R_4\) instance, proving \(R_4^*\). These are fixed
finite primitive-rule constructions, so the rules are derived, not
merely semantically admissible.
\end{proof}
\end{annexitem}

\begin{annexitem}{annex:p18.qev.admissible-rules}{thm:p18-qev-admissible-rules}{All eleven admissible transformations}
\begin{proof}
By soundness, provable premise sequents evoke their questions.
For each rule we prove evocation of the conclusion. Its antecedent
is finite and its question is well formed by the displayed provisos;
completeness then gives a proof of the conclusion sequent.

\emph{Right complementary removal (Fact 15).} If an \(X\)-valuation
made every common \(A_i\) false, coverage of the two premise lists
would force both \(B\) and \(\neg B\) true. This is impossible.
Thus the common list is covered. No common \(A_i\) is entailed,
by either premise. The condition \(n>1\) ensures a question remains.

\emph{Left complementary cases (Fact 16).} Any \(X\)-valuation
satisfies \(B\) or \(\neg B\), so coverage of the appropriate
premise gives a true answer. For each \(i\), a countervaluation
to entailment of \(A_i\) from \(X\cup\{B\}\), supplied by
the first premise, is also an \(X\)-countervaluation. Hence no
individual answer follows from \(X\).

\emph{Left equivalence (Fact 17).} If \(B\leftrightarrow C\)
is a tautology, \(X\cup\{B\}\) and \(X\cup\{C\}\)
have exactly the same satisfying valuations. Both coverage and each
individual countervaluation are unchanged.

\emph{Reverse \(R_3\) (Fact 18).} This is the forward-to-right
direction of the conditionalization equivalence already proved.
The common prefix preserves question syntax.

\emph{Right implication transformation (Fact 19).} Suppose the
question with answers \(A_i\to B\) is evoked by \(X\).
An \(X\cup\{\neg B\}\)-valuation must satisfy some
\(A_i\to B\); since \(B\) is false, it satisfies \(\neg A_i\).
For each \(i\), unsettledness of \(A_i\to B\) gives a
valuation satisfying \(X\), making \(A_i\) true and \(B\)
false. This is a countervaluation to entailment of \(\neg A_i\)
from \(X\cup\{\neg B\}\). The conclusion is therefore evoked.

\emph{Left \(\beta\)-combination (Fact 20).} A valuation satisfying
\(X\cup\{\beta\}\) satisfies \(\beta_1\) or \(\beta_2\)
by the defining truth equivalence. Coverage of the corresponding
premise provides a true answer. Each countervaluation for \(A_i\)
under \(X\cup\{\beta_1\}\) is also one under
\(X\cup\{\beta\}\), since \(\beta_1\models\beta\).
Thus all individual non-entailments hold.

\emph{Both left \(\alpha\)-transformations (Facts 21 and 22).}
The premise sets \(X\cup\{\alpha\}\) and
\(X\cup\{\alpha_1,\alpha_2\}\) have exactly the same
valuations. Coverage and each non-entailment agree in both directions.
This proves both contraction to \(\alpha\) and expansion into
its two components, including cases where a component already belongs
to the set antecedent.

\emph{Right \(\beta\)-expansion (Fact 23).} Coverage of the old
list gives a common answer or \(\beta\). In the latter case
\(\beta_1\) or \(\beta_2\) is true, so the expanded list
is covered. Neither new component is entailed: entailment of
\(\beta_j\) would imply entailment of \(\beta\), contrary
to the premise. Common answers remain unentailed. The syntactic
provisos ensure the expanded list is a question. This argument works
also with exactly one common answer, giving the source's additional
small-question instance.

\emph{Right \(\alpha\)-combination (Fact 24).} If an
\(X\)-valuation makes every common answer false, coverage of the
two premise questions forces both \(\alpha_1\) and
\(\alpha_2\) true; hence \(\alpha\) is true. Conversely,
if \(X\models\alpha\), then \(X\models\alpha_1\),
contrary to unsettledness of the first premise. Common answers remain
unentailed, again by a premise. This proves evocation and also works
with one common answer. No assumption that the literal expression
\(\alpha_1\land\alpha_2\) differs from the common answers
is needed: admissibility uses the stated formula \(\alpha\)
and its own syntactic proviso.

\emph{Export (Fact 25).} Let \(X\cup\{B\}\) evoke the
original question. For any \(X\)-valuation, either \(B\) is
false, making the added answer \(\neg B\) true, or \(B\)
is true and a common answer is true by premise coverage. For each
common \(A_i\), its countervaluation under \(X\cup\{B\}\)
also witnesses \(X\not\models A_i\). Such a countervaluation
exists for at least one answer and makes \(B\) true, so
\(X\not\models\neg B\). The new list is therefore covered
and wholly unsettled. Its additional syntactic condition is exactly
the proviso on the sheet.

Finally, all three rows of the \(\alpha/\beta\) table have the
asserted equivalences. The conjunction row is its truth clause; the
negated-disjunction and negated-conjunction rows follow by negating
their binary truth conditions; the implication rows follow because
\(B\to C\) is false exactly when \(B\) is true and \(C\)
false. This checks the classical premises used in the last six
transformations and completes all eleven cases.
\end{proof}
\end{annexitem}

\begin{annexitem}{annex:p18.qev.source-examples}{thm:p18-qev-source-examples}{The five primitive-rule proofs}
\begin{proof}
Each row below records a sequent in a proof and the primitive rule
that yields it. Formula letters introduced within a case abbreviate
the explicitly fixed displayed formulas, not additional assumptions.

\emph{First example.} Put \(S=p\lor\neg p\). Since \(S\)
is tautological, \(S\to p\) is equivalent to \(p\), and
\(S\to\neg p\) is equivalent to \(\neg p\). The proof is
\[
\begin{array}{ll}
\vdash?\{p,\neg p\}&\mathrm{Ax}\\
\vdash?\{p,S\to\neg p\}&R_2\\
\vdash?\{S\to\neg p,p\}&R_4\\
\vdash?\{S\to\neg p,S\to p\}&R_2\\
S\vdash?\{\neg p,p\}&R_3\\
S\vdash?\{p,\neg p\}&R_4.
\end{array}
\]

\emph{Second example.} Put \(S=p\lor q\). The classical
truth functions give
\(S\to p\leftrightarrow(\neg q\lor p)\) and
\(S\to q\leftrightarrow(\neg p\lor q)\): if \(p\) is
true the first implication and its right-hand side are true; if
\(p\) is false they both have value \(\neg q\). The second
equivalence is symmetric. Thus
\[
\begin{array}{ll}
\vdash?\{\neg q\lor p,\neg p\lor q\}&\mathrm{Ax}\\
\vdash?\{\neg q\lor p,S\to q\}&R_2\\
\vdash?\{S\to q,\neg q\lor p\}&R_4\\
\vdash?\{S\to q,S\to p\}&R_2\\
S\vdash?\{q,p\}&R_3\\
S\vdash?\{p,q\}&R_4.
\end{array}
\]

\emph{Third example.} Let
\(H=p\to(q\lor r)\),
\(F_q=H\to(p\to q)\), and
\(F_r=H\to(p\to r)\). Expanding the classical implication
shows that \(F_q\) is equivalent to \(\neg r\lor\neg p\lor q\):
if \(p\) is false both are true; with \(p\) true the remaining
condition is \((q\lor r)\to q\), equivalent to
\(\neg r\lor q\). Similarly \(F_r\) is equivalent to
\(\neg q\lor\neg p\lor r\). Hence
\[
\begin{array}{ll}
\vdash?\{\neg r\lor\neg p\lor q,\neg q\lor\neg p\lor r\}&\mathrm{Ax}\\
\vdash?\{\neg r\lor\neg p\lor q,F_r\}&R_2\\
\vdash?\{F_r,\neg r\lor\neg p\lor q\}&R_4\\
\vdash?\{F_r,F_q\}&R_2\\
\vdash?\{F_q,F_r\}&R_4\\
H\vdash?\{p\to q,p\to r\}&R_3\\
H,p\vdash?\{q,r\}&R_3.
\end{array}
\]

\emph{Fourth example.} Let
\(H=(p\land q)\to r\),
\(G_p=H\to(\neg r\to\neg p)\), and
\(G_q=H\to(\neg r\to\neg q)\).
The formula \(G_p\) is equivalent to \(q\lor r\lor\neg p\).
If \(r\) is true or \(p\) is false, both are true; in the
remaining case \(r=0,p=1\), the implication has value \(q\).
The corresponding argument gives
\(G_q\leftrightarrow(p\lor r\lor\neg q)\). Therefore
\[
\begin{array}{ll}
\vdash?\{q\lor r\lor\neg p,p\lor r\lor\neg q\}&\mathrm{Ax}\\
\vdash?\{q\lor r\lor\neg p,G_q\}&R_2\\
\vdash?\{G_q,q\lor r\lor\neg p\}&R_4\\
\vdash?\{G_q,G_p\}&R_2\\
\vdash?\{G_p,G_q\}&R_4\\
H\vdash?\{\neg r\to\neg p,\neg r\to\neg q\}&R_3\\
H,\neg r\vdash?\{\neg p,\neg q\}&R_3.
\end{array}
\]

\emph{Fifth example.} Both
\(\vdash?\{p,q,\neg p\}\) and
\(\vdash?\{p,q,\neg q\}\) are axioms. Apply \(R_1\)
with common list \(p,q\) to obtain
\(\vdash?\{p,q,\neg p\land\neg q\}\). By the truth
conditions of disjunction and negation,
\(\neg p\land\neg q\leftrightarrow\neg(p\lor q)\)
is a tautology. One \(R_2\) application gives the required final
question.

Every axiom here satisfies the complementary-literal criterion:
no single displayed clause contains a complementary pair, but the
whole answer list does. Distinctness of \(p,q,r\) also makes all
initial clauses distinct. During the replacements a conditional
expression cannot coincide with a remaining literal or disjunction;
the two completed conditional answers differ in their final variables
or negated variables. The last example's conjunction and negated
disjunction differ syntactically from \(p,q\). Thus all \(R_1,R_2\)
provisos hold, not just their semantic equivalences.
\end{proof}
\end{annexitem}

\begin{annexitem}{annex:p18.qev.decision}{thm:p18-qev-decision}{Finite decision and nonmonotonicity of evocation}
\begin{proof}
There are finitely many input variables and exactly \(2^r\) Boolean
assignments to them. Each formula has a finite parse tree, so recursive
evaluation terminates for every assignment. Any such assignment extends
to a valuation of the full language by assigning zero to all remaining
variables. Conversely, every full valuation restricts to an enumerated
assignment, and the truth values of all input formulas are preserved.

Thus the coverage test succeeds exactly when every full valuation of
the premises satisfies an answer. Recording a false instance of each
answer succeeds exactly when each answer has a countervaluation under
the premises. These are exactly the two evocation conditions. If the
premises are inconsistent, there are no recorded witnesses and the
method correctly rejects. Adequacy converts the same result into
decision of PMCE provability.

For failure of unrestricted weakening, take \(X=\{p\lor q\}\)
and \(Q=?\{p,q\}\). The two countervaluations described in the
body, together with coverage, prove \(E(X,Q)\). Adding \(p\)
makes that answer entailed, so \(E(X\cup\{p\},Q)\) fails.
Adequacy gives a provable first sequent and a nonprovable second one.
Ordinary classical consequence remains monotone because restricting
the satisfying valuations of an antecedent cannot introduce a new
countervaluation to a conclusion it already entails. Evocation fails
to be monotone precisely because it also requires countervaluations
to each individual answer.
\end{proof}
\end{annexitem}
